\documentclass[preprint,12pt]{elsarticle}

\usepackage{amsmath,amssymb,amsthm,booktabs,graphicx}
\usepackage{microtype}
\usepackage{hyperref}
\hypersetup{hidelinks}
\usepackage{siunitx}
\usepackage{enumitem}
\usepackage{placeins}

\journal{Journal of Computational and Applied Mathematics}

\newtheorem{theorem}{Theorem}
\newtheorem{proposition}[theorem]{Proposition}
\newtheorem{lemma}[theorem]{Lemma}
\newtheorem{corollary}[theorem]{Corollary}
\theoremstyle{remark}
\newtheorem{remark}[theorem]{Remark}

\begin{document}

\begin{frontmatter}

\title{Model-Aware Positive Stieltjes Quadrature:
From Structural Minimax to Gauss--Legendre}

\author[inst1]{Tran Minh Duc\corref{cor1}}
\cortext[cor1]{Corresponding author}
\ead{ductranminh2107@gmail.com}

\affiliation[inst1]{
  organization={FPT University},
  addressline={Lot E2a-7, D1 Street, Saigon Hi-Tech Park},
  city={Ho Chi Minh City},
  country={Vietnam}
}

\begin{abstract}
Generalized Gaussian and rationally exact quadrature already exploit
non-polynomial structure, while rational minimax approximation under
interpolation constraints is an established approximation problem. We
consider positive symmetric Stieltjes-response rules that minimize
relative error uniformly over a continuum of Lorentzian widths while
matching a variable prefix of polynomial moments. Within this class,
\(n\) positive node pairs and moment depth \(k\) give the feasible range
\(0\leq k\leq2n\), admissible dimension \(2n-k\), and the
Gauss--Legendre rule at \(k=2n\). For \(k<2n\), a
difference-numerator degree bound yields a \(2n-k+1\)-point
alternation certificate. High-precision hierarchies for \(n=2,3,4\)
quantify the trade-off between structural accuracy and polynomial
protection. For a smooth residual with Chebyshev coefficients
\(|a_j|\leq C\rho^{-j}\), an exact coefficient-space norm gives the
model-aware risk \(G_k(\rho)+\tau S_k\). Arbitrary-precision ball
arithmetic encloses all nine structural maxima and the Chebyshev
defects of the rounded \(n=4\) operators. Adaptive interval comparison
then certifies a unique selector on 99.5956\% of the
tested \((\rho,\log_{10}\tau)\)-coordinate area, leaving
0.4044\% unresolved near transition curves. A
multi-start sensitivity study and a shifted-Lorentzian stress test
support numerical reliability and local, rather than universal,
robustness.
\end{abstract}

\begin{keyword}
numerical quadrature \sep Stieltjes functions \sep rational minimax
approximation \sep moment constraints \sep Gauss--Legendre quadrature
\sep interval arithmetic
\end{keyword}

\end{frontmatter}

\section{Introduction}

Polynomial exactness is a central organizing principle of numerical
quadrature, but it is not the only way to spend a fixed node budget.
When the dominant difficulty is a nearby pole, quadrature may instead
incorporate rational functions chosen to match that structure
\cite{gautschi1993}. Generalized Gaussian quadrature extends positive
exactness to Chebyshev systems and is naturally described through
moment-space geometry \cite{huybrechs2017}. Positive interpolatory
formulas and optimization-based moment-matching constructions provide
additional general frameworks \cite{vandebos2019,glaubitz2020,
keshavarzzadeh2018}.

A second neighboring direction begins with rational approximation.
Modern algorithms compute rational minimax approximants through adaptive
barycentric representations and Remez-type iterations
\cite{filip2018}. Interpolation-constrained rational minimax
approximation, including higher-order constraints, is also established
\cite{zhang2025}. For Markov and Stieltjes functions, relative rational
approximation has a substantial theory \cite{beckermann2021}. Most
directly, quadrature formulas can be derived from rational
approximations of Cauchy transforms, with poles and residues becoming
nodes and weights and Gauss quadrature connected to Pad\'e approximation
at infinity \cite{horning2025}.

These results rule out several broad novelty claims: pole-aware
quadrature, positive exact quadrature, rational minimax,
equioscillation, and constrained rational approximation are not new
separately. The narrower problem studied here is to combine them in a
model-aware design. At fixed node-pair count, we require a positive
symmetric Stieltjes response, match a variable prefix of polynomial
moments, and minimize relative error uniformly over a continuum of
Lorentzian widths. The resulting family spans a structural minimax rule
and the Gauss--Legendre endpoint. The practical question is then not
only how to compute the rules, but how to select the moment depth when a
smooth uncertain residual is superposed on the structural component.

Within this stated class, the paper makes the following contributions:
\begin{enumerate}[label=(\roman*)]
\item it formulates the positive variable-moment continuum-minimax
problem and characterizes its feasible range, dimension, and
Gauss--Legendre endpoint;
\item it gives a short class-specific degree/zero-count argument leading
to a \(2n-k+1\)-point alternation certificate for every
nonzero-dimensional level;
\item it computes the complete path for \(n=2,3,4\), including the
upper-half levels between the earlier lower-half designs and the
Gaussian endpoint;
\item it derives an exact mixed structural/smooth risk and uses it to
select the moment depth across a two-parameter model domain;
\item it supplies a conservative interval-certified decision map, a shifted-pole
scope test, and multi-start reconstruction from perturbed initializations.
\end{enumerate}

The feasible-range theorem and the Gaussian endpoint are important
structural components of the framework, but are not claimed as
stand-alone inventions of moment or generalized Gaussian theory. A
targeted literature search did not identify a direct source combining
positivity, continuum Lorentzian minimax, variable moment depth, an exact
mixed selector, and a interval-certified parameter partition. This is not an
exhaustive proof of novelty; specialist prior-art review remains
necessary. The claims are restricted to the symmetric positive class
and do not imply universal superiority over Gaussian, adaptive, or
problem-specific quadrature.

\section{Problem formulation}

Let
\[
I[f]=\int_{-1}^{1} f(x)\,dx,
\qquad
L_b(x)=\frac{1}{x^2+b^2},
\qquad
b\in B=[0.08,0.12].
\]
Set
\[
S=\{b^2:b\in B\}=[0.08^2,0.12^2].
\]
With \(s=b^2\) and \(y=x^2\),
\[
F(s)=I[L_{\sqrt{s}}]
=\int_0^1\frac{y^{-1/2}}{s+y}\,dy
=\frac{2}{\sqrt{s}}\arctan\frac{1}{\sqrt{s}},
\qquad s\in S.
\]

For a fixed integer \(n\geq1\), define the ordered positive parameter
domain
\[
\mathcal P_n
=
\left\{
(y,w)\in\mathbb R^{2n}:
0<y_1<\cdots<y_n<1,\quad w_j>0
\right\}.
\]
Each \((y,w)\in\mathcal P_n\) determines the symmetric rule
\[
Q_{y,w}[f]
=
\sum_{j=1}^{n}w_j
\left\{
f(-\sqrt{y_j})+f(\sqrt{y_j})
\right\}.
\]
Its response on the structural family is
\[
r_{y,w}(s)
=
Q_{y,w}[L_{\sqrt{s}}]
=
2\sum_{j=1}^{n}\frac{w_j}{s+y_j}.
\]

For an integer \(k\geq0\), let
\[
m^{(k)}
=
\left(\frac{1}{2\ell+1}\right)_{\ell=0}^{k-1}
\]
and define the moment map
\[
M_k(y,w)
=
\left(
\sum_{j=1}^{n}w_jy_j^\ell
\right)_{\ell=0}^{k-1}.
\]
For \(k=0\), both vectors are empty. The admissible parameter set and
the corresponding rule class are
\[
\Theta_{n,k}
=
\left\{
(y,w)\in\mathcal P_n:
M_k(y,w)=m^{(k)}
\right\},
\qquad
\mathcal A_{n,k}
=
\left\{
Q_{y,w}:(y,w)\in\Theta_{n,k}
\right\}.
\]

\begin{lemma}[Equivalent interpretations of the moment constraints]
Let \(k\geq1\) and \((y,w)\in\mathcal P_n\). The following statements
are equivalent:
\begin{enumerate}[label=(\alph*)]
\item \((y,w)\in\Theta_{n,k}\);
\item \(Q_{y,w}\) is exact for every polynomial of degree at most
\(2k-1\);
\item
\[
r_{y,w}(s)-F(s)=O(s^{-k-1})
\qquad
\text{as }s\to+\infty.
\]
\end{enumerate}
\end{lemma}

\begin{proof}
Symmetry gives exactness for odd polynomials. For even monomials,
\[
Q_{y,w}[x^{2\ell}]
=
2\sum_{j=1}^{n}w_jy_j^\ell,
\qquad
I[x^{2\ell}]
=
\frac{2}{2\ell+1}.
\]
This proves the equivalence of (a) and (b).

For every fixed \(m\geq1\), uniformly for \(y\in[0,1]\),
\[
\frac{1}{s+y}
=
\sum_{\ell=0}^{m-1}
(-1)^\ell y^\ell s^{-\ell-1}
+
O(s^{-m-1}).
\]
Consequently,
\[
F(s)
=
2\sum_{\ell=0}^{m-1}
\frac{(-1)^\ell}{2\ell+1}s^{-\ell-1}
+
O(s^{-m-1}),
\]
whereas
\[
r_{y,w}(s)
=
2\sum_{\ell=0}^{m-1}
(-1)^\ell
\left(\sum_{j=1}^{n}w_jy_j^\ell\right)
s^{-\ell-1}
+
O(s^{-m-1}).
\]
Taking \(m=k\) proves that (a) implies (c). Conversely, if (c) holds,
uniqueness of the coefficients in these finite asymptotic expansions
forces equality of the first \(k\) moments, proving (c) implies (a).
\end{proof}

\begin{lemma}[Uniqueness of the representation]
If \((y,w),(\widehat y,\widehat w)\in\mathcal P_n\) and
\[
r_{y,w}(s)=r_{\widehat y,\widehat w}(s)
\]
on a nonempty open real interval, then
\[
(y,w)=(\widehat y,\widehat w).
\]
\end{lemma}

\begin{proof}
The identity theorem makes the two rational functions identical. Their
simple poles are \(-y_j\) and \(-\widehat y_j\), and the corresponding
residues are \(2w_j\) and \(2\widehat w_j\). Equality of poles and
residues, together with the imposed ordering, gives equality of the
parameter vectors.
\end{proof}

\section{Feasible-range characterization}

\begin{theorem}[Feasible range, regular dimension, and Gaussian endpoint]
For \(n\geq1\), the following statements hold:
\begin{enumerate}[label=(\roman*)]
\item \(\Theta_{n,k}\neq\varnothing\) for every integer
\(0\leq k\leq2n\);
\item for \(0\leq k\leq2n\), \(\Theta_{n,k}\) is an embedded smooth
submanifold of the open set \(\mathcal P_n\), of dimension
\[
\dim\Theta_{n,k}=2n-k;
\]
\item \(\Theta_{n,k}=\varnothing\) for every integer \(k\geq2n+1\);
\item \(\Theta_{n,2n}\) is a singleton, and its associated symmetric
rule is the classical \(2n\)-point Gauss--Legendre rule.
\end{enumerate}
\end{theorem}

\begin{proof}
Let \(d\mu(y)=y^{-1/2}\,dy\) on \([0,1]\). By the standard Gaussian
quadrature theorem for a positive measure with infinite support, there
exists a unique \(n\)-node rule with positive weights that is exact for
all polynomials of degree at most \(2n-1\). Write its nodes and weights
as \(y_j\) and \(2w_j\). Then
\[
\sum_{j=1}^{n}w_jy_j^\ell
=
\frac{1}{2}\int_0^1y^\ell\,d\mu(y)
=
\frac{1}{2\ell+1},
\qquad
0\leq\ell\leq2n-1.
\]
Thus \(\Theta_{n,2n}\neq\varnothing\), and every lower moment prefix is
also feasible. The case \(k=0\) is simply
\(\Theta_{n,0}=\mathcal P_n\).

For \(1\leq k\leq2n\), suppose a linear combination of the rows of
\(DM_k(y,w)\) vanishes, and define
\[
p(t)=\sum_{\ell=0}^{k-1}c_\ell t^\ell.
\]
The columns corresponding to \(w_j\) give \(p(y_j)=0\). The columns
corresponding to \(y_j\) give \(w_jp'(y_j)=0\), and positivity gives
\(p'(y_j)=0\). Hence every \(y_j\) is a double root of \(p\). Since
\(\deg p\leq k-1\leq2n-1\), this is possible only when \(p\equiv0\).
Therefore \(DM_k\) has row rank \(k\) at every point of
\(\mathcal P_n\). The regular-level-set theorem gives the stated
embedded-manifold dimension.

If \(k\geq2n+1\), the constraints include exactness for the transformed
polynomial
\[
g(y)=\prod_{j=1}^{n}(y-y_j)^2,
\]
which has degree \(2n\), is nonnegative, and is not identically zero.
The discrete rule gives
\[
2\sum_{j=1}^{n}w_jg(y_j)=0,
\]
whereas
\[
\int_0^1g(y)\,d\mu(y)>0,
\]
a contradiction.

It remains to identify the endpoint uniquely. Let
\[
\pi_n(y)=\prod_{j=1}^{n}(y-y_j)
\]
for any element of \(\Theta_{n,2n}\). For every polynomial \(q\) of
degree at most \(n-1\), exactness through degree \(2n-1\) gives
\[
\int_0^1\pi_n(y)q(y)\,d\mu(y)
=
2\sum_{j=1}^{n}w_j\pi_n(y_j)q(y_j)
=
0.
\]
Thus \(\pi_n\) is the unique monic degree-\(n\) orthogonal polynomial
for \(\mu\), so the nodes are uniquely determined. Exactness for
degrees \(0,\ldots,n-1\) then determines the weights uniquely through a
nonsingular Vandermonde system. Under \(y=x^2\), the resulting
symmetric \(2n\)-node rule is exact for all polynomials of degree at
most \(4n-1\), and therefore is the \(2n\)-point Gauss--Legendre rule.
\end{proof}

\begin{remark}
The sets \(\Theta_{n,k}\) are nested as moment constraints are added.
The selected minimax rules form a discrete hierarchy indexed by \(k\);
no connectedness of the level sets, and no continuous path of
minimizers in \(k\), is asserted.
\end{remark}

\section{Minimax characterization}

For \((y,w)\in\mathcal P_n\), write
\[
q_y(s)=\prod_{j=1}^{n}(s+y_j),
\qquad
r_{y,w}(s)=\frac{p_{y,w}(s)}{q_y(s)},
\]
where \(\deg p_{y,w}\leq n-1\).

\begin{proposition}[Difference-numerator degree]
Let \(Q_1,Q_2\in\mathcal A_{n,k}\) with \(0\leq k\leq2n\). Then either
their responses are identical or
\[
r_{Q_1}(s)-r_{Q_2}(s)
=
\frac{H(s)}{q_1(s)q_2(s)},
\qquad
\deg H\leq2n-k-1.
\]
For \(k=2n\), the second alternative is impossible.
\end{proposition}

\begin{proof}
The unreduced numerator
\[
H=p_1q_2-p_2q_1
\]
has degree at most \(2n-1\). For \(k=0\), both responses are
\(O(s^{-1})\), so their difference is \(O(s^{-1})\). For \(k\geq1\),
the moment lemma gives
\[
r_{Q_1}(s)-r_{Q_2}(s)=O(s^{-k-1}).
\]
Since \(q_1q_2\) is monic of degree \(2n\), a nonzero polynomial
numerator must satisfy
\[
\deg H-2n\leq-k-1,
\]
which is the claimed bound.
\end{proof}

Define
\[
\varepsilon_Q(s)=\frac{r_Q(s)}{F(s)}-1,
\qquad
\|\varepsilon_Q\|_{\infty,S}
=
\max_{s\in S}|\varepsilon_Q(s)|.
\]

\begin{theorem}[Class-specific alternation certificate]
Let \(0\leq k<2n\) and let \(Q^\star\in\mathcal A_{n,k}\). Suppose
there are
\[
m=2n-k+1
\]
distinct ordered points
\[
s_1<\cdots<s_m
\quad\text{in }S
\]
and a number \(E>0\) such that
\[
\varepsilon_{Q^\star}(s_i)=(-1)^i\sigma E,
\qquad
\sigma\in\{-1,1\},
\]
and
\[
E=\|\varepsilon_{Q^\star}\|_{\infty,S}.
\]
Then no rule in \(\mathcal A_{n,k}\) has strictly smaller uniform
relative error on \(S\).
\end{theorem}

\begin{proof}
Assume that \(Q\in\mathcal A_{n,k}\) satisfies
\[
\|\varepsilon_Q\|_{\infty,S}<E.
\]
At every \(s_i\), the difference
\[
\varepsilon_{Q^\star}(s_i)-\varepsilon_Q(s_i)
\]
has the sign of \(\varepsilon_{Q^\star}(s_i)\). These signs alternate,
so continuity gives at least \(m-1=2n-k\) distinct zeros in the open
subintervals \((s_i,s_{i+1})\). Since \(F(s)>0\) on \(S\), these are
zeros of \(r_{Q^\star}-r_Q\). The difference-numerator proposition
allows at most \(2n-k-1\) distinct zeros unless the responses are
identical. In the identical-response case, uniqueness of the ordered
positive representation gives \(Q=Q^\star\), which is incompatible
with strict improvement. This contradiction proves the result.
\end{proof}

\begin{remark}
The theorem is a sufficient global-minimax certificate within
\(\mathcal A_{n,k}\). It does not prove necessity of alternation,
uniqueness among equal-norm minimizers, or existence of an exact
equioscillating parameter vector from floating-point residuals. The
endpoint \(k=2n\) is zero-dimensional and is fixed by its moments.
\end{remark}


\section{Numerical realization of the model-aware hierarchy}

The theorem above is conditional on exact moment
constraints and exact equioscillation. Numerically, we solve the
corresponding nonlinear systems at high precision and report residual
audits. We denote the ideal mathematical rule by \(Q_{n,k}^\star\) and
its published decimal operator by \(\widetilde Q_{n,k}\). The table and
figures use the computed approximations; they are numerical evidence
for the conditional theorem, not a computer-assisted existence proof.

The nonlinear systems impose the moment equations, alternating
equal-magnitude errors, and stationarity at internal extrema. Lower-half rules were obtained first, and the upper-half
problems were continued toward the Gaussian endpoint. All reported
operators have ordered nodes and positive weights.

\begin{table}[ht]
\centering
\caption{Complete \(n=4\) hierarchy. The endpoint \(k=8\) is
zero-dimensional; therefore the alternation column is not applicable.}
\label{tab:full-hierarchy}
\small
\begin{tabular}{cccccc}
\toprule
\(k\) & Dimension & Exactness & Alt. points &
Worst relative error & Min. \(w_j\)\\
\midrule
0 & 8 & -- & 9 & 1.490623e-10 & 0.039147 \\
1 & 7 & 1 & 8 & 2.460626e-09 & 0.044528 \\
2 & 6 & 3 & 7 & 4.048641e-08 & 0.051670 \\
3 & 5 & 5 & 6 & 6.653546e-07 & 0.061646 \\
4 & 4 & 7 & 5 & 1.092819e-05 & 0.076595 \\
5 & 3 & 9 & 4 & 1.794333e-04 & 0.101407 \\
6 & 2 & 11 & 3 & 2.943219e-03 & 0.149011 \\
7 & 1 & 13 & 2 & 4.597982e-02 & 0.126194 \\
8 & 0 & 15 & -- & 4.300499e-01 & 0.101229 \\
\bottomrule
\end{tabular}
\end{table}

\begin{figure}[ht]
\centering
\includegraphics[width=0.84\linewidth]{figures/figure1_full_hierarchy.pdf}
\caption{Worst centered relative error for the full feasible
hierarchies \(0\leq k\leq2n\), \(n=2,3,4\). The final point of each
curve is the corresponding Gauss--Legendre rule.}
\label{fig:full-hierarchy}
\end{figure}

For \(n=4\), the error grows from \(1.49\times10^{-10}\) at \(k=0\) to
\(4.30\times10^{-1}\) at the Gauss endpoint. The upper-half values
\(k=5,6,7\) continue the approximately constant finite-range price per
additional moment before the error enters an order-one regime. This
pattern is empirical and is not claimed as an asymptotic theorem.

\begin{figure}[ht]
\centering
\includegraphics[width=0.82\linewidth]{figures/figure2_alternation.pdf}
\caption{Representative normalized computed alternation patterns from the lower
endpoint (\(k=0\)), midpoint (\(k=4\)), and upper half (\(k=7\)) of the
\(n=4\) hierarchy.}
\label{fig:alternation}
\end{figure}

\FloatBarrier

\section{Exact coefficient-space rule selection}

For \(C\geq0\) and \(\rho>1\), define
\[
\mathcal H(C,\rho)
=
\left\{
h:
h(x)=\sum_{j=0}^{\infty}a_jT_j(x),
\quad
|a_j|\leq C\rho^{-j}
\right\}.
\]
Because \(|T_j(x)|\leq1\) on \([-1,1]\) and
\(\sum_j\rho^{-j}<\infty\), every series in this definition converges
absolutely and uniformly and therefore represents a continuous
function.

For a fixed bounded quadrature functional \(Q\), define
\[
\Delta_j(Q)=I[T_j]-Q[T_j]
\]
and
\[
G_Q(\rho)
=
\sum_{j=0}^{\infty}\rho^{-j}|\Delta_j(Q)|.
\]
The series converges because
\[
|\Delta_j(Q)|
\leq
\|I-Q\|_{C[-1,1]\to\mathbb R}.
\]

\begin{theorem}[Exact coefficient-box norm]
For every fixed bounded quadrature functional \(Q\),
\[
\sup_{h\in\mathcal H(C,\rho)}
|I[h]-Q[h]|
=
C G_Q(\rho).
\]
\end{theorem}

\begin{proof}
Uniform convergence permits termwise application of the bounded
functional \(I-Q\), so
\[
(I-Q)h
=
\sum_{j=0}^{\infty}a_j\Delta_j(Q).
\]
The coefficient bounds give
\[
|(I-Q)h|
\leq
C\sum_{j=0}^{\infty}
\rho^{-j}|\Delta_j(Q)|.
\]
Equality is attained by choosing
\[
a_j
=
C\rho^{-j}\operatorname{sign}\Delta_j(Q),
\]
with any value in \([-C\rho^{-j},C\rho^{-j}]\) when
\(\Delta_j(Q)=0\). The resulting series belongs to
\(\mathcal H(C,\rho)\).
\end{proof}

For \(c_{\max}\geq0\), define the parameterized mixed class
\[
\mathcal F(C,\rho,c_{\max})
=
\left\{
h+cL_b:
h\in\mathcal H(C,\rho),\
|c|\leq c_{\max},\
b\in B
\right\}.
\]
For a fixed rule, let
\[
S_Q
=
\max_{b\in B}|I[L_b]-Q[L_b]|.
\]
The maximum exists because \(B\) is compact and the error is
continuous in \(b\).

\begin{theorem}[Exact mixed-class risk]
For every fixed bounded quadrature functional \(Q\),
\[
\sup_{f\in\mathcal F(C,\rho,c_{\max})}
|I[f]-Q[f]|
=
C G_Q(\rho)+c_{\max}S_Q.
\]
\end{theorem}

\begin{proof}
The upper bound follows from the triangle inequality and the preceding
theorem. Choose \(b_\star\in B\) attaining \(S_Q\). Select
\(h_\star\in\mathcal H(C,\rho)\) such that
\[
(I-Q)h_\star=C G_Q(\rho),
\]
and choose
\[
c_\star
=
c_{\max}
\operatorname{sign}
\left(I[L_{b_\star}]-Q[L_{b_\star}]\right).
\]
Then both error contributions have the same sign and their absolute
values add, proving attainability.
\end{proof}

Let \(Q_0,\ldots,Q_{2n}\) be the finite hierarchy of candidate rules.
When \(C>0\), put \(\tau=c_{\max}/C\). The exact worst-case selector
within this candidate set is
\[
k^\star(\rho,\tau)
\in
\arg\min_{0\leq k\leq2n}
\left\{
G_{Q_k}(\rho)+\tau S_{Q_k}
\right\}.
\]

\begin{remark}
This selector is exact for the stated coefficient box and structural
family, but only among the precomputed candidates
\(\{Q_0,\ldots,Q_{2n}\}\). It does not assert global optimality over all
quadrature rules. The rules were constructed by a relative structural
minimax objective, whereas \(S_Q\) is the absolute structural error
appropriate to the additive mixed class.
\end{remark}


\begin{table}[ht]
\centering
\caption{Full \(n=4\) structural and smooth-risk trade-off at
\(\rho=1.5\).}
\label{tab:risk-components}
\small
\begin{tabular}{ccc}
\toprule
\(k\) & Structural screening margin \(S_k\) & \(G_k(1.5)\)\\
\midrule
0 & 5.55627e-09 & 9.11242e-01 \\
1 & 9.17189e-08 & 3.95465e-01 \\
2 & 1.50911e-06 & 2.04657e-01 \\
3 & 2.48007e-05 & 9.59856e-02 \\
4 & 4.07341e-04 & 4.49242e-02 \\
5 & 6.68824e-03 & 2.14018e-02 \\
6 & 1.09706e-01 & 1.17665e-02 \\
7 & 1.71386e+00 & 6.62192e-03 \\
8 & 1.60298e+01 & 3.42291e-03 \\
\bottomrule
\end{tabular}
\end{table}

\begin{figure}[ht]
\centering
\includegraphics[width=0.78\linewidth]{figures/figure3_risk_tradeoff.pdf}
\caption{Each added moment decreases the smooth-risk norm but increases
the Lorentzian structural error. Labels indicate the moment depth
\(k\).}
\label{fig:risk-tradeoff}
\end{figure}

\section{Computer-assisted interval certification for the rounded rules}

The exact risk theorems apply to each fixed operator. For the published
decimal operators \(\widetilde Q_k\), we treat every displayed decimal
node and weight as an exact decimal input enclosed by ball arithmetic.
All transcendental and algebraic evaluations in the structural stage
use 80-decimal-digit Python-FLint/Arb balls
\cite{johansson2017}.

For each \(k=0,\ldots,8\), let
\[
e_k(b)
=
I[L_b]-\widetilde Q_k[L_b].
\]
The interval \([0.08,0.12]\) is adaptively bisected. On each subinterval
\(X\), two enclosures are computed: the natural interval extension
\(e_k(X)\), and the mean-value enclosure
\[
e_k(m)+e'_k(X)(X-m),
\]
where \(m\) is the midpoint. The smaller valid absolute upper bound is
used in a branch-and-bound search. Point evaluations provide lower
witnesses. This produces rigorous enclosures
\[
S_k
=
\max_{b\in[0.08,0.12]}|e_k(b)|
\in[\underline S_k,\overline S_k].
\]
All nine largest lower witnesses occur at \(b=0.08\); the maximum
relative enclosure width is
\[
1.437e-10.
\]

For the smooth component, each Chebyshev defect
\[
\Delta_j(\widetilde Q_k)
=
I[T_j]-\widetilde Q_k[T_j]
\]
is enclosed by Arb through degree \(J=500\). Since
\[
|\Delta_j(\widetilde Q_k)|
\leq
\|I-\widetilde Q_k\|
\leq
2+2\sum_iw_i,
\]
the remaining tail is bounded by
\[
\sum_{j>J}\rho^{-j}|\Delta_j|
\leq
\left(2+2\sum_iw_i\right)
\frac{\rho^{-(J+1)}}{1-\rho^{-1}}.
\]
The nonnegative finite sums are evaluated by outward-rounded Horner
recurrences. On a cell
\[
[\rho_-,\rho_+]\times[t_-,t_+],
\qquad t=\log_{10}\tau,
\]
monotonicity gives rigorous lower and upper risk bounds from
\(\rho_+,\rho_-\) and \(10^{t_-},10^{t_+}\), respectively. A cell is
assigned to \(k\) only when the upper bound for \(k\) is strictly below
the lower bound for every competitor.

Adaptive binary subdivision to depth
\(d=24\) visits
344,615 cells and certifies a
unique winner on
\[
99.595636\%
\]
of the rectangular coordinate area
\[
(\rho,\log_{10}\tau)
\in[1.05,5]\times[-12,12].
\]
The remaining
\[
0.404364\%
\]
is unassigned and lies near transition curves. The percentage is an
area fraction in these coordinates, not a probability statement.

\begin{table}[ht]
\centering
\caption{Fraction of the interval-certified coordinate area assigned
to each moment depth.}
\label{tab:selection-fractions}
\small
\begin{tabular}{cc}
\toprule
\(k\) & Fraction of certified area\\
\midrule
0 & 21.993\% \\
1 & 9.179\% \\
2 & 8.320\% \\
3 & 8.493\% \\
4 & 8.536\% \\
5 & 8.572\% \\
6 & 8.312\% \\
7 & 7.054\% \\
8 & 19.540\% \\
\bottomrule
\end{tabular}
\end{table}

\begin{figure}[ht]
\centering
\includegraphics[width=0.84\linewidth]{figures/figure4_certified_selector.pdf}
\caption{Interval-certified full-hierarchy selector. Blank regions
remain unresolved because candidate risk intervals overlap.}
\label{fig:selector}
\end{figure}

The certificate applies to the rounded operators, coefficient box,
structural interval, and coordinate domain stated above. It does not
prove existence or uniqueness of exact equioscillating parameter
vectors for the ideal nonlinear minimax problems.

\FloatBarrier

\section{External validity and reconstruction}

\subsection{Small pole-location mismatch}

Without retraining, the centered rules were tested on
\[
L_{a,b}(x)=\frac{1}{(x-a)^2+b^2},
\qquad
|a|\leq0.03,
\quad
0.08\leq b\leq0.12.
\]

\begin{table}[ht]
\centering
\caption{Worst relative errors on the shifted-Lorentzian stress-test
box.}
\label{tab:shifted}
\small
\begin{tabular}{lcccc}
\toprule
Method & Evaluations & Worst error & Worst \(|a|\) & Worst \(b\)\\
\midrule
PMC--k0 & 8 & 4.522336e-06 & 0.0300 & 0.0800 \\
PMC--k4 & 8 & 1.025564e-03 & 0.0300 & 0.0845 \\
Composite--Simpson--9 & 9 & 7.426762e-02 & 0.0300 & 0.1200 \\
Gauss--Legendre--8 & 8 & 4.300499e-01 & 0.0000 & 0.0800 \\
\bottomrule
\end{tabular}
\end{table}

\begin{figure}[ht]
\centering
\includegraphics[width=0.82\linewidth]{figures/figure5_shifted_stress.pdf}
\caption{Stress test under small pole-location mismatch. The comparison
supports local structural robustness, not shifted-family minimaxity or
universal superiority.}
\label{fig:shifted}
\end{figure}

The \(k=0\) rule has worst relative error
\(4.52\times10^{-6}\), whereas the Gauss endpoint has worst error
\(4.30\times10^{-1}\) on the same box. The centered error is nevertheless
inflated by more than four orders of magnitude, so the correct
interpretation is local robustness rather than shift invariance.

\subsection{Multi-start reconstruction}

Using the same equioscillation equations and solver implementation, the six upper-half problems
\[
(2,3),\ (3,4),\ (3,5),\ (4,5),\ (4,6),\ (4,7)
\]
were reconstructed from twelve perturbed starts each. A run was accepted
only when nodes were ordered, weights were positive, moment residuals
were below \(10^{-10}\), and scaled equioscillation residuals were below
\(10^{-8}\).

\begin{table}[ht]
\centering
\caption{Multi-start reconstruction of the upper-half
rules. Error differences are relative to the published minimax level.}
\label{tab:multistart}
\scriptsize
\begin{tabular}{ccccc}
\toprule
\((n,k)\) & Accepted starts & Error difference &
Max. node difference & Max. weight difference\\
\midrule
$(2,3)$ & 12/12 & 2.64e-16 & 5.55e-17 & 5.55e-17 \\
$(3,4)$ & 10/12 & 6.78e-14 & 4.44e-16 & 5.00e-16 \\
$(3,5)$ & 12/12 & 5.31e-14 & 1.29e-14 & 1.32e-14 \\
$(4,5)$ & 8/12 & 2.38e-12 & 4.44e-15 & 4.27e-15 \\
$(4,6)$ & 11/12 & 9.22e-14 & 1.03e-14 & 1.05e-14 \\
$(4,7)$ & 12/12 & 5.11e-13 & 8.95e-14 & 6.43e-14 \\
\bottomrule
\end{tabular}
\end{table}

All accepted solutions in each problem formed a single error-level
cluster. No accepted start found a better rule by more than
\(10^{-6}\) relative, and the maximum best-rule error difference was
\(2.38\times10^{-12}\). This is a sensitivity test with perturbed initializations. It reduces the risk of an isolated local numerical artifact, but it is neither an independent implementation nor a proof of global uniqueness.

\section{Relation to prior work and novelty scope}

\subsection{Structure-aware and generalized quadrature}

Gautschi's rational Gauss-type rules establish the principle that a
quadrature formula can integrate rational functions selected to match
important external poles \cite{gautschi1993}. The present work therefore
does not introduce pole-aware quadrature. Its structural objective is
different: it minimizes relative error over a continuum of Lorentzian
widths instead of enforcing exactness at a finite set of prescribed
poles.

Generalized Gaussian quadrature for Chebyshev systems supplies positive
interior-node rules exact for non-polynomial spaces, and complete
Chebyshev systems support stepwise increases in exactness
\cite{huybrechs2017}. Positive exact formulas and positive-node
geometry are also well developed \cite{vandebos2019,glaubitz2020}.
Accordingly, neither positivity nor an ordered hierarchy of exactness is
claimed as new. Here the moment depth indexes a family selected by a
continuum Stieltjes minimax criterion at fixed node-pair count.

Optimization-based nested quadrature shows that nonlinear
moment-matching systems can be solved flexibly
\cite{keshavarzzadeh2018}. Our nonlinear solver and multi-start study
are reproducibility mechanisms, not the principal conceptual
contribution.

\subsection{Rational approximation and minimax constraints}

Rational minimax approximation and equioscillation are classical;
adaptive barycentric algorithms make high-order problems computationally
accessible \cite{filip2018}. Interpolation-constrained rational minimax
approximation already includes prescribed function values and
higher-order derivative conditions \cite{zhang2025}. After the change
of variable \(z=1/s\), matching Laurent coefficients at infinity is
closely related to confluent interpolation at \(z=0\). Thus the novelty
claim cannot be ``minimax plus constraints.'' The alternation theorem in
this paper is instead a concise specialization to positive
Stieltjes-response quadrature, where the shared coefficients reduce the
difference-numerator degree.

Markov-function rational approximation provides relevant relative-error
context \cite{beckermann2021}. Horning and Trefethen make the link from
approximation to quadrature explicit by deriving quadrature
from rational approximations of Cauchy transforms; their framework also
connects Gauss quadrature with Pad\'e approximation at infinity
\cite{horning2025}. This is the closest conceptual comparator. The
distinction here is the simultaneous enforcement of positive residues,
a variable polynomial-moment prefix, a continuum relative-minimax
objective, and an exact mixed-class selector.

\subsection{Worst-case rule selection}

Worst-case analysis for analytic function classes is an established
part of quadrature theory \cite{trefethen2022,goda2024}, and recent
positive-weight kernel quadrature gives a different
convex-geometric worst-case framework \cite{hayakawa2026}. The
coefficient-box identity used here is elementary duality. Its role is
to make the smooth residual commensurate with the Stieltjes structural
error, producing
\[
G_k(\rho)+\tau S_k.
\]
The paper's practically distinctive output is therefore not a new norm
identity in isolation, but a selector and interval-certified decision map over a
finite family of positive structural rules.

\subsection{Defensible novelty statement}

The targeted audit found no direct source combining all of the
following: positive symmetric Stieltjes residues, continuum Lorentzian
relative minimax, variable moment depth at fixed node-pair count, an
exact mixed smooth/structural selector, and a validated parameter
partition. The defensible claim is this combined model-aware framework.
The paper does not claim to invent positivity, moment matching,
structure-aware quadrature, constrained rational minimax,
equioscillation, or the Gauss/Pad\'e connection separately.

\section{Limitations}

The theoretical and numerical study is one-dimensional and symmetric.
The principal structural family is the centered Lorentzian kernel. The
shifted-family experiment is a stress test and is not itself
minimax-optimized.

The alternation theorem is conditional: the interval selector
certificate does not by itself prove that exact equioscillating
upper-half rules exist near the reported floating-point parameters.
Global uniqueness among equal-norm minimizers is not proved. The
multi-start study uses the same nonlinear equations and solver
implementation, so it is not independent replication.

The interval-certified selector covers 99.5956\% of the
chosen \((\rho,\log_{10}\tau)\)-coordinate rectangle and leaves
0.4044\% unassigned. The coverage depends on the
specified coefficient box, Lorentzian interval, candidate rules,
coordinate domain, truncation degree \(J=500\), and subdivision depth.
It is not a probability statement and does not imply optimality over
quadrature rules outside the finite hierarchy.

The empirical near-geometric price per moment is not an asymptotic
theorem. The literature audit is targeted rather than exhaustive, and
equivalent formulations may exist under moment-space, generalized
Gaussian, constrained Markov, or interpolatory rational terminology.
The manuscript therefore avoids ``first'' and ``unprecedented'' claims.

\section{Conclusion}

Within the stated positive symmetric Stieltjes class, the variable
moment depth organizes a path
\[
k=0
\quad\longrightarrow\quad
k=2n
\]
from a structural minimax rule to Gauss--Legendre. The
feasible-range characterization and difference-degree argument support
the numerical construction, while the exact mixed risk turns the family
into a model-aware selector. Complete computations for \(n=2,3,4\),
computer-assisted interval certification of the decision map, a shifted-pole stress test, and
multi-start reconstruction quantify both the promise and the limits of
the approach.

The main contribution is not any ingredient separately. It is the
combination of positive Stieltjes realization, continuum structural
minimax, variable polynomial protection, and exact rule selection in one
computable framework. The next scientific step is independent
specialist comparison with the closest generalized-quadrature and
constrained-rational literature, followed by extensions beyond symmetry
and the centered kernel.

\section*{Data and code availability}

The code, published rules, generated data, and reproduction scripts are
available in the
\href{https://github.com/TurkeyOp/partial-moment-stieltjes-quadrature}{public GitHub repository}.
The Version 2.3 interval-certification and full-selector materials are included in the
accompanying research archive and should be merged into the public
repository before journal submission.

\section*{Funding}

This research did not receive any specific grant from funding agencies
in the public, commercial, or not-for-profit sectors.

\section*{Declaration of competing interest}

The author declares that he has no known competing financial interests
or personal relationships that could have appeared to influence the
work reported in this paper.

\section*{Declaration of generative AI and AI-assisted technologies in
the manuscript preparation process}

During the preparation of this work, the author used OpenAI ChatGPT to
assist with literature organization, code prototyping, manuscript
structuring, and language revision. The author reviewed and verified
the mathematical derivations, numerical results, references,
interpretations, and final wording and takes full responsibility for
the content of the article.

\begin{thebibliography}{14}
\providecommand{\url}[1]{\texttt{#1}}

\bibitem{gautschi1993}
W. Gautschi, Gauss-type quadrature rules for rational functions,
arXiv:math/9307223 (1993).

\bibitem{huybrechs2017}
D. Huybrechs, On the computation of Gaussian quadrature rules for
Chebyshev sets of linearly independent functions,
arXiv:1710.11244 (2017).

\bibitem{filip2018}
S.-I. Filip, Y. Nakatsukasa, L. N. Trefethen, B. Beckermann,
Rational minimax approximation via adaptive barycentric
representations, arXiv:1705.10132 (2017).

\bibitem{zhang2025}
L.-H. Zhang, Y.-N. Zhang, b-d-Lawson: a method for the interpolation
constrained rational minimax approximation,
arXiv:2502.10665 (2025).

\bibitem{horning2025}
A. Horning, L. N. Trefethen, Quadrature formulas from rational
approximations, arXiv:2507.14971 (2025).

\bibitem{vandebos2019}
L. M. M. van den Bos, B. Sanderse, A geometric approach for the
addition of nodes to an interpolatory quadrature rule with positive
weights, arXiv:1902.07477 (2019).

\bibitem{glaubitz2020}
J. Glaubitz, Constructing positive interpolatory cubature formulas,
arXiv:2009.11981 (2020).

\bibitem{keshavarzzadeh2018}
V. Keshavarzzadeh, R. M. Kirby, A. Narayan, Generation of nested
quadrature rules for generic weight functions via numerical
optimization: Application to sparse grids, arXiv:1808.03707 (2018).

\bibitem{beckermann2021}
B. Beckermann, J. Bisch, R. Luce, On the rational approximation of
Markov functions, with applications to the computation of Markov
functions of Toeplitz matrices, arXiv:2106.05098 (2021).

\bibitem{trefethen2022}
L. N. Trefethen, Exactness of quadrature formulas,
arXiv:2101.09501 (2021).

\bibitem{goda2024}
T. Goda, Y. Kazashi, K. Tanaka, How sharp are error bounds? Lower
bounds on quadrature worst-case errors for analytic functions,
arXiv:2401.07196 (2024).

\bibitem{hayakawa2026}
S. Hayakawa, Convex-geometric error bounds for positive-weight kernel
quadrature, arXiv:2605.05705 (2026).

\bibitem{nailwal2024}
R. Nailwal, A. Zalar, Gaussian quadratures with prescribed nodes via
moment theory, arXiv:2412.20849 (2024).


\bibitem{johansson2017}
F. Johansson, Arb: efficient arbitrary-precision midpoint-radius
interval arithmetic, IEEE Transactions on Computers 66 (2017)
1281--1292; arXiv:1611.02831.

\end{thebibliography}

\end{document}
